\appendix
\section{Logical-space and distance proofs}
\label{app:slab-proof}

All vector spaces in these proofs are over $\mathbb F_2$. Logical Pauli
classes are taken modulo stabilizers and irrelevant phases. Let $g_X,g_Z$
denote the numbers of listed lattice stars and faces, and $r_X^{\rm rel}$,
$r_Z^{\rm rel}$ the dimensions of their relation spaces. These relation
dimensions are distinct from the numbers of disjoint ports in the main text.

\subsection{Commutation and relation maps}
The junction rules in Sec.~\ref{sec:construction-details} preserve
commutation except at a crossing of an $A_i$ sheet and a $C_j$ sheet.
For example, a data star and an adjacent $C$ face share both the data
$z$-edge and the added $c_y$ edge. An $A$ star and an adjacent data face
share both the added $d_z$ and an $a_x$ edge. At an $A_i$--$C_j$ crossing,
each common input qubit contributes one unmatched overlap. The wall
equations add $\sigma_{y'}+\tau_y$ to the overlap of the star at $y'$ and
the face on $[y,y+1]$, for $y'=y,y+1$. Successive equations determine
the variables along a path, with the endpoint equations specified in
Sec.~\ref{sec:construction-details}. Their sum is the parity of
$\operatorname{supp}X_i\cap\operatorname{supp}Z_j$. This is zero by
$H_XH_Z^T=0$, so the endpoint conditions agree and the solution is unique.
An empty overlap has no variables; for a one-vertex overlap only the two
endpoint conditions on the same variable remain. Thus all checks commute.

The sheet counts, assigning shared boundary columns to data sheets, are
\begin{align}
 n_D&=L_x(L_z-1)+(L_x+1)L_z,\nonumber\\
 g_{X,D}&=(L_x+1)(L_z-1),\quad g_{Z,D}=L_xL_z,\nonumber\\
 n_C&=w(L_z-1)+(w-1)L_z,\nonumber\\
 g_{X,C}&=(w-1)(L_z-1),\quad g_{Z,C}=wL_z,\nonumber\\
 n_A&=L_x(w+1)+(L_x+1)w,\nonumber\\
 g_{X,A}&=(L_x+1)(w+1),\quad g_{Z,A}=L_xw.
\end{align}
Here $w$ is the sheet's width in lattice edges. Data sheets contribute
$+1$ and either check sheet contributes $-1$ to the Euler count, so
\begin{equation}
 n-g_X-g_Z=n_{\rm in}-m_X-m_Z.
 \label{eq:rank-euler}
\end{equation}
The junction and wall additions change supports without adding generators
or qubits. Redundant input rows cannot be discarded in this count.

For an $X$ relation, assign a binary selection coefficient to every star.
Each edge owned by $A_i$ belongs to its two endpoint $A_i$ stars and no
other star, so the selection is constant, say $\alpha_i$, on that sheet.
On the data column $x=0$, a data $z$-edge belongs to the data stars at its
endpoints when present, and to an $A_i$ star at a corresponding incidence.
Summing the vanishing edge coefficients along the column, where there are
no stars at $z=0,L_z$, gives
\begin{equation}
 \sum_{i:q\in X_i}\alpha_i=0\quad\text{for every }q,
 \qquad \alpha^TH_X=0.
\end{equation}
The coefficients on the data stars are determined from the bottom
boundary; data $x$-edges make them constant along each row. Coefficients
on the remaining $C$ stars are then determined from their bottom
boundary, since $c_z(y,0)$ meets just one such star. Thus the map from a
lattice $X$ relation to $\alpha$ is injective.

For a $Z$ relation, the edges owned by $C_j$ imply a constant face
coefficient $\gamma_j$ on that sheet. Data $x$-edges imply that the data
face coefficient is independent of $z$; write it as $b_q(x)$.
The data $z$-edges give
\begin{align}
 b_q(0)&=b_q(L_x-1)=0,\nonumber\\
 b_q(x-1)+b_q(x)&=\sum_{j:X_j=x,\ q\in Z_j}\gamma_j.
\end{align}
Summing yields $\gamma^TH_Z=0$. The data coefficients are determined by
$\gamma$, and the $A$-face coefficients are determined from $x=0$, where
each $a_y$ edge meets one face. Hence this relation map is also injective.
Consequently,
\begin{align}
 k&=n-g_X-g_Z+r_X^{\rm rel}+r_Z^{\rm rel}\nonumber\\
 &\le n_{\rm in}-m_X-m_Z+(m_X-\rank H_X)\nonumber\\
 &\quad +(m_Z-\rank H_Z)\nonumber\\
 &=K.
 \label{eq:rank-upper}
\end{align}

\subsection{Slab pairing, stabilizers, and translation}
The $X$ slab for $x\in\ker H_Z$ contains all data $z$-edges at height
$z_0$ on selected data sheets. On $C_j$, the segment after incidence $t$
has coefficient $p_t=\sum_{h\le t}x_{q_h}$. More precisely it contains
the $C_j$-owned $z$-edges with $Y_{q_t}<y\le Y_{q_{t+1}}$, excluding
the shared last column. Local overlaps with faces are the endpoint and
middle relations of Eqs.~\eqref{eq:boundary} and~\eqref{eq:prefix}.
They all vanish. Walls add only $a_y$ edges to these faces and hence
do not change this calculation.

For $z\in\ker H_X$, the dual slab at $x_0$ contains every data $z$-edge
at that $x$ coordinate on selected data sheets and the $A_i$-owned
$y$-edges with $Y_{q_t}\le y<Y_{q_{t+1}}$, weighted by the analogous
prefix parity. It commutes with every star. Opposite slabs intersect on
exactly one data edge for each input coordinate selected by both; their
other supports are sheet owned and disjoint. Their pairing is therefore
\begin{equation}
 \langle\widetilde X(x;z_0),\widetilde Z(z;x_0)\rangle=x\cdot z.
 \label{eq:slab-pair-proof}
\end{equation}
Choose input logical bases with pairing the identity. Their lifts give
$K$ independent logical pairs and thus $k\ge K$. Together with
Eq.~\eqref{eq:rank-upper}, this proves $k=K$.

This also closes the two identifications needed in the main text without
assuming them in the rank argument. If $x$ is an input $X$ stabilizer,
its lift pairs trivially with every lifted $Z$ logical. Those lifted
logicals span the full opposite logical space because $k=K$, so the lift
is a lattice stabilizer. The dual statement follows identically.
The difference of two heights of an $X$ lift has zero pairing with every
$Z$ lift by Eq.~\eqref{eq:slab-pair-proof}, so it too is a stabilizer.
The same holds for translating $Z$ planes. Hence the lifts induce an
isomorphism of input and lattice logical spaces at every allowed position.

\subsection{The per-plane projection}
Fix a nontrivial lattice $X$ logical $O$ and a plane $x=x_0$. Decompose
the dual slab into disjoint pieces: one vertical data column per input
qubit and one $A_i$ backbone segment per interval between successive
incidences. Let $a_q$ and $b_{i,t}$ be the parities of the intersection of
$O$ with these pieces. Define the input vector
\begin{equation}
 \pi_{x_0}(O)=\sum_q a_q e_q+
       \sum_{i,t}b_{i,t}\sum_{h\le t}e_{q_h^{(X_i)}}.
 \label{eq:plane-projection}
\end{equation}
For every $z\in\ker H_X$, its dot product with $z$ is the commutator
of $O$ with the $Z$ lift of $z$ in this plane. Taking $z$ to be an input
$Z$ stabilizer proves $\pi_{x_0}(O)\in\ker H_Z$. Taking a logical $z$
conjugate to $O$ proves that the projected class is nontrivial. This
argument holds in every plane because the lifts represent the same class.

Replacing a prefix by its complementary suffix adds the input $X_i$
stabilizer. Choose the shorter representative of each prefix; it has
weight at most $\lfloor w_X/2\rfloor$. Each parity-one piece contains
at least one occupied qubit of $O$, and distinct pieces in a plane do
not share qubits. Thus a representative of $\pi_{x_0}(O)$ has weight at
most $\eta_X$ times the number of occupied qubits of $O$ in the plane.
Its weight is at least $d_X^{\rm in}$. Summing over the $L_x+1$ disjoint
planes proves Eq.~\eqref{eq:distance}. Exchanging the two constructions
and summing over the $L_z$ disjoint layers proves
Eq.~\eqref{eq:distance-dual}. Complete slabs have $O(N^2)$ qubits for
balanced bounded-degree inputs: there are $O(N)$ data pieces of length
$O(N)$ and $O(N)$ bounded-weight checks with total backbone length
$O(N^2)$. This supplies the matching upper bound for every pure logical
class and, by taking a product of two slabs, for every mixed Pauli class.

\section{Grid expansion and unequal-distance protection}
\label{app:expansion}

\subsection{A discrete relative-expansion bound}
Let $P_a$ be the path with $a$ vertices. Side lengths in this appendix
count vertices, not edges. Let $G_0=P_a\square P_b$ with $2\le a\le b$,
and let $h(G_0)$ be its edge Cheeger constant. We first prove
\begin{equation}
 h(G_0)\ge 1/b.
\end{equation}
For $|S|\le ab/2$, regard the $a$ rows as having length $b$. If there is
both a full and an empty row, every column has a vertical boundary edge,
so $|\delta S|\ge b\ge |S|/b$. If no row is full, every nonempty row
has a horizontal boundary edge, giving $|\delta S|\ge |S|/b$.
If a full row exists but no row is empty, all columns meet $S$ and at
most $|S|/a\le b/2$ columns are full. At least $b/2$ columns therefore
have a vertical boundary edge, again giving $|\delta S|\ge |S|/b$.
Complementation handles larger sets.

Now take $G=G_0\square P_\ell$ and let $U$ occupy at most $t$ of its
parallel $G_0$ layers, with $|U|\ge2$. For any cut set $S$, denote its
layer slices by $S_j$. Fix one layer $j_0$, and put
\begin{equation}
 \mu=\min\{|S_{j_0}|,|V(G_0)\setminus S_{j_0}|\},\qquad
 m=\max_j|S_j\mathbin\triangle S_{j_0}|.
\end{equation}
The vertical boundary has at least $m$ edges: along a path of layers
to a maximizing slice, every changed vertex crosses a boundary.
Each horizontal slice has minority size at least $\max\{0,\mu-m\}$.
Writing $c=\min\{1,\ell h(G_0)\}$, the total boundary obeys
\begin{equation}
 |\delta S|\ge m+c\max\{0,\mu-m\}\ge c\max\{\mu,m\}.
\end{equation}
Choose the minority side of the reference slice; in each port layer it
contains at most $\mu+m$ vertices. Hence
\begin{equation}
 \min\{|S\cap U|,|U\setminus S|\}
       \le t(\mu+m)\le2t\max\{\mu,m\}.
\end{equation}
If $\max(a,b,\ell)/\min(a,b,\ell)\le A$, then
$c\ge1/A$ and $\beta(G,U)\ge1/(2At)$, proving
Eq.~\eqref{eq:grid-bound}. Any coordinate orientation can be used.
Ports need not fill their layers. Constant lattice refinement separates
the finitely many local port labels and changes $A,t$ only by constants.
Cases with fewer than two ports are excluded from this relative-expansion
definition and treated directly.

\subsection{Parallel edges and the cycle space}
Replace every edge $e$ by $m_0$ parallel copies and attach all copies to
the same vertex stars. Use one designated copy for each original pairing
path and each elementary square cycle. In addition, measure
$Z_{e,1}Z_{e,r}$ for $2\le r\le m_0$. These checks span the entire
cycle space of the multigraph: a cycle can first be moved to designated
copies using the two-edge cycles and then reduced using the rectangular
grid's square cycles. Thus no additional cycle logicals are introduced.
Every cut contains all $m_0$ copies of each original cut edge, so
\begin{equation}
 \beta(G^{(m_0)},U)=m_0\beta(G,U).
\end{equation}
Taking $m_0=\lceil2At\rceil$ makes this at least one. At fixed $m_0$,
star weights, cycle weights, and check incidences remain bounded. The
copies can be placed in a fixed refinement cell near the original edge;
bounded congestion then permits a constant-depth local extraction schedule
for each static check catalog. Uniform properties across an entire adaptive
history still require Assumption~\ref{ass:uniform}.

\subsection{The protected quotient for unequal blocks}
Here use the dual convention and gauge
$M=\overline Z_1\overline Z_2=Z(D_1)Z(D_2)$ on $C_1\otimes C_2$.
Let $C_2$ have distance at least $d_2$. Include all graph cycle checks.
Let $\mathcal Q_1$ be generated by deformed images of the logical
operators on $C_1$ that commute with $\overline Z_1$. An $X$ image is
completed by edges with boundary equal to its port incidence; a $Z$
image needs no completion. Different completions differ by a cycle
stabilizer. Because $M$ is measured, $\overline Z_2$ and
$\overline Z_1$ represent the same class in $\mathcal Q_1$.

\begin{lemma}[Static protected quotient]
The minimum weight of a deformed logical class outside $\mathcal Q_1$
is at least $\min\{1,\beta(G,U)\}d_2$.
\end{lemma}
\begin{proof}
For an $X$ normalizer, restriction to the original blocks is in their
original $X$ normalizers; the edge boundary equation also requires
commutation with $M$. If its $C_2$ restriction is a stabilizer, remove
its deformed lift. The remaining $C_1$ logical commutes with
$\overline Z_1$, and differences of edge completions are cycles, so
the class lies in $\mathcal Q_1$. Otherwise its $C_2$ data restriction
already has weight at least $d_2$.

For a $Z$ normalizer, the ancilla support is a cut $\delta S$, since
it is orthogonal to every cycle. Multiplying by vertex stars on either
side removes the cut and gives two original logicals differing by $M$.
If the $C_2$ part of either belongs to
$\langle S_Z(C_2),\overline Z_2\rangle$, the class is in
$\mathcal Q_1$. Outside this subgroup both reductions have nontrivial
$C_2$ parts, each of weight at least $d_2$. Put $a$ equal to the
original data weight, $m=\min\{|S\cap U|,|U\setminus S|\}$, and
$b_0=\min\{1,\beta\}$. The smaller-side reduction has weight at
most $a+m$ and at least $d_2$, whereas
\begin{equation}
 a+|\delta S|\ge a+b_0m\ge b_0(a+m)\ge b_0d_2.
\end{equation}
Finally, $\mathcal Q_1$ is generated by pure-type classes in this CSS
setting. A mixed class outside it has a pure component outside it, of
weight no larger than the mixed operator, giving the same bound.
\end{proof}

The lemma also permits $C_1$ to be a bare qubit, when its commuting
logical subgroup is generated by $Z_1$. During teleportation the parity
and final $X_1$ outcomes update only the traveling slot. This algebraic
statement explains the exceptional resource degree of freedom. It does
not by itself prove that every inherited syndrome at a noisy open
boundary is harmless to spectators; that requirement is stated separately
in Assumption~\ref{ass:uniform}.

\section{A relative phenomenological slice bound}
\label{app:slice}

Consider a rectangle with at least one undeformed entry and exit round
and $R$ deformed rounds gauging $P=X(D)$. A fault is a single-qubit
Pauli in a gap before a round or a flip of one check record. Edge
preparation and readout flips are represented at their first and last
gaps. Subtract the fault-free propagation of the same incoming
$(s,\lambda,\rho)$ and fix the same ideal random gauge outcomes.
The relative input is zero; require zero relative syndrome and records
at the right cut. The first and last cut records are included as
boundary data, not silently replaced by perfect physical measurements.

Write $E_t$ for the accumulated data-and-edge Pauli by round $t$.
For a check $c$, with commutator $\sigma_c(E_t)$ and record fault
$m_{c,t}$, its relative record is
\begin{equation}
 \Delta r_{c,t}=\sigma_c(E_t)+m_{c,t}.
\end{equation}
Use consecutive-round detectors, preparation and readout relations, and
the local entry/exit relations between a $Z$ check $s$ and its lift
$\widetilde s=sZ(\mu_s)$. The measured cycles generate the full cycle
space. We prove Eq.~\eqref{eq:rect-distance} for a detector-silent pattern.

The CSS sectors can be treated separately; a $Y$ mechanism contributes
to both sectors but its count is not doubled when bounding either part.
For the $X$ sector, consecutive and entry/exit equations propagate zero
relative $Z$-check records from the cuts through the rectangle. Let $z_e$
be the parity of $X$ faults on edge $e$, which is its relative final
$Z$ readout. Exit equations give
\begin{equation}
 \sum_{e\in\mu_s}z_e=0,\qquad
 \sum_{e\in c}z_e=0\quad\text{for every cycle }c.
\end{equation}
Thus $z=\delta S$. The final data error $E_{\rm d}$ commutes with the
original $Z$ checks by the output-syndrome condition. The single-slice
operator $O=E_{\rm d}X(\delta S)$ commutes with the deformed checks.
Multiplying its vertex stars over $S$ leaves
$L=E_{\rm d}X(f^{-1}(S\cap U))$, exactly the relative data frame
after edge-readout correction. Complementing $S$ changes $L$ by $P$,
which is trivial after the measurement. A nonzero surviving effect means
that $O$ is a nontrivial deformed logical, so $|O|\ge\widetilde d_X$.
Every occupied data or edge timeline in $O$ contains at least one original
$X$-sector fault. Therefore this sector needs at least
$\widetilde d_X$ faults for a nontrivial effect.

For the $Z$ sector, original $X$-check records are chained from the zero
input, and each vertex-star relative record is a constant $b_u$ across
the deformed rounds. The reported parity flip is $o=\sum_u b_u$.
The output data error $E_{\rm d}$ is an original $Z$ normalizer. If its
class is nontrivial, its data support already costs at least $d_Z$
faults. Otherwise a nonzero effect requires $o=1$. If every deformed
round contains an $X$-check record fault, there are at least $R$ faults.
If some round $t^*$ has none, the data part $Z(B)$ of $E_{t^*}$ commutes
with every original $X$ check. Each edge occurs in two vertex stars,
whereas each port occurs in one, so
\begin{equation}
 |B\cap D|=\sum_u\sigma_{A_u}(E_{t^*})=\sum_u b_u=1.
\end{equation}
Hence $Z(B)$ anticommutes with $P$ and is a nontrivial original logical
of weight at least $d_Z$. Its occupied data timelines cost at least that
many faults. This sector therefore needs $\min\{d_Z,R\}$ faults.
At least one sector has a nonzero effect, proving the claimed minimum.
With no deformation, the same output argument gives the idle-memory
distance $\min\{d_X,d_Z\}$. The $X/Z$-exchanged version applies to
$Z$ gauging. This is a lower bound for the stated rectangle and boundary
conditions; exact attainment in each particular circuit requires its own
witness or distance calculation.

\section{Deterministic resource factories}
\label{app:factories}

\subsection{Codes, orientations, and recovery}
For $r=3$ or $4$, let $n_r=2^r-1$ and label coordinates by the distinct
nonzero vectors $v_i\in\mathbb F_2^r$ in binary order. Let $L_r$ be the
row span of the $r\times n_r$ matrix with columns $v_i$, and define a
CSS code with
\begin{equation}
 S_X=L_r,\qquad S_Z=\langle L_r,\mathbf1\rangle^\perp.
\end{equation}
Here $S_X,S_Z$ specify binary supports of stabilizer generators. The
code encodes one qubit, with the fixed computational-basis convention
\begin{equation}
 |b_L\rangle=2^{-r/2}\sum_{x\in L_r}|x+b\mathbf1\rangle,
 \qquad b\in\{0,1\}.
 \label{eq:resource-codewords}
\end{equation}
For $r=3$ this is the Steane code; for $r=4$ it is the punctured
Reed--Muller code. All nonzero words of $L_r$ have weight $2^{r-1}$.
Consequently $(S^\dagger)^{\otimes7}$ has phase one on $|0_L\rangle$
and phase $i$ on $|1_L\rangle$, while
$(T^\dagger)^{\otimes15}$ has phase one and $e^{i\pi/4}$ respectively.
They implement logical $S$ and $T$ in this convention.

Prepare $|+_L\rangle$, apply the indicated transversal rotations by
coordinatewise resource teleportation, measure the $X$-check syndrome,
recover, and decode the logical output. The CSS encoder and decoder are
fixed by Eq.~\eqref{eq:resource-codewords}; Gaussian elimination realizes
them with CNOTs, $X/Z$ preparations, and $X/Z$ measurements. They require
no extra non-Clifford candidate slots. For the negative-angle rotations
one may use $X|Y\rangle$ or $X|A\rangle$, equal up to phase to the
corresponding conjugate resource. Each conditional $S^\dagger$ correction
after a $T^\dagger$ injection is assigned one additional $Y$ slot on
the same coordinate.

The recovery table is completely specified by
\begin{equation}
 R(0)=I,\qquad R(v_i)=Z_i.
 \label{eq:factory-recovery}
\end{equation}
The columns $v_i$ enumerate every nonzero syndrome once. Thus recovery
is defined on every outcome and the factory always returns one output.
In ideal operation all $Z$-check syndromes are zero; measuring them can
also be included without changing this phase-error recovery rule.

\subsection{Arbitrary defective candidates}
Purify a candidate as $|0\rangle|\phi_0\rangle+
|1\rangle|\phi_1\rangle$. In any fixed phase-teleportation branch, after
the prescribed Pauli correction, the data-to-output map is diagonal in
the computational basis. Relative to the ideal diagonal rotation it is
therefore $I\otimes E_0+Z\otimes E_1$. Tracing out the environment gives
conditional Kraus operators in $\operatorname{span}\{I,Z\}$. This proof
also covers mixed and correlated candidates by purification. A defective
correction candidate remains diagonal on the same coordinate.

If at most one coordinate $i$ is defective, the encoded block before
recovery has error space
$I\otimes E_0+Z_i\otimes E_1$. The two syndrome outcomes are $0$ and
$v_i$; Eq.~\eqref{eq:factory-recovery} corrects both. After correction
the retained logical state is ideal and factors from the environment.
This is a statement about correction of an error space, not about a
particular stochastic mixture of $I$ and $Z$.

For a resource family with defect-set decomposition
\begin{equation}
 \rho=\sum_B\pi(B)
   \left(\bigotimes_{i\notin B}|R_i\rangle\langle R_i|\right)
   \otimes\sigma_{B,E},\quad
 \sum_{B\supseteq J}\pi(B)\le q^{|J|},
\end{equation}
a bad output requires at least two defective slots, in fact at distinct
coordinates. For any $u$ disjoint output runs, choose a pair in each;
their $2u$ slots are distinct. The joint defect bound and union bound give
\begin{equation}
 \Pr(\text{all }u\text{ outputs bad})
 \le \binom{30}{2}^{u}q^{2u}.
\end{equation}
This proves Eq.~\eqref{eq:distill} for the entire output family. The
constant 435 is a conservative slot-pair count. It is consistent with
the 21 and 105 distinct-coordinate pairs for the two codes and does not
claim the cubic error suppression of postselected distillation.

Let $y_\ell,a_\ell$ count raw candidates needed for one $Y$ or $A$
output at level $\ell$. With preassigned correction slots,
\begin{equation}
 y_\ell=7y_{\ell-1},\qquad
 a_\ell=15a_{\ell-1}+15y_{\ell-1},\qquad y_0=a_0=1.
\end{equation}
Since $7/15<1$, both are $O(15^\ell)$. The number 30 in the error
bound therefore does not change the production exponent to
$\log_2 30$. These recurrences count ideal logical factory operations
up to a fixed factor; physical faults in those operations are addressed
by the scheduled-circuit hypothesis.

\section{Conditional threshold and resource accounting}
\label{app:conditional}

This appendix proves the architecture implication under
Assumption~\ref{ass:uniform}. It does not use finite certificates as
proofs of that assumption. Let $z$ bound the degree of the fixed union
fault graph. The number of connected sets of $m$ vertices containing a
specified anchor is at most $C_z^m$ for a size-independent constant
$C_z$. One can see this by encoding a rooted spanning tree and its
bounded-degree neighbor choices; overcounting suffices.

If at least a fraction $\theta$ of the vertices in such a set are actual
faults, the local-stochastic condition and a union bound over their
subsets give probability at most $2^m p^{\theta m}$. Subsets with
incompatible labels at one location have probability zero, so they do
not weaken the bound. Hence a large-cluster event over $V$ possible
anchors is bounded by a geometric sum of
$V(2C_zp^\theta)^m$. Choosing a fixed threshold makes this sum
exponentially small in its minimum size. Item 4 of the assumption
supplies the required size and density properties for committed decoding;
they are not consequences of bounded graph degree alone.

For clarity, the half-density statement in a fixed nonadaptive linear
problem follows from minimum cardinality. If $e$ is the actual error and
$\widehat e$ a minimum solution with the same syndrome, then
$e+\widehat e$ is detector silent. Each of its connected components is
silent. Replacing the decoder's part on one component by the actual part
remains a valid solution, so the actual part is at least as large as the
decoder part there. At least half the residual component is therefore
actual. Open commitment boundaries, mutually exclusive location labels,
and changes of adaptive branch require the corresponding uniform witness
statement rather than an unqualified extension of this argument.

\subsection{Joint courier error estimate}
Take a fixed-factor sequence
$\lambda_-d_j\le d_{j+1}\le\lambda_+d_j$ with $1<\lambda_-\le\lambda_+$,
and $R_j=\Theta(d_j)$. Set $u=C\sqrt p<1$. The anchor and half-density
bounds imply a one-resource charging estimate
\begin{equation}
 q_{\rm inj}\le C_0u+
 C_1\sum_{j=0}^{J}d_j^{5/2}u^{\gamma d_j}
 +C_1f^{5/2}u^{\gamma d_J}.
 \label{eq:courier-sum-final}
\end{equation}
The first term is the raw stage. Choose the constant starting distance
with $\gamma d_0\ge1$. For $u$ below a fixed constant, the geometric
growth of $d_j$ makes the sum $O(u)$ uniformly in $J$. Choose
\begin{equation}
 d_J\ge d_{\rm base}
      +C_{\rm chain}\frac{\log f}{\log(1/p)}
\end{equation}
with a sufficiently large fixed coefficient. Since
$\log(1/u)/\log(1/p)$ is bounded below for sufficiently small $p$,
the last term is also $O(\sqrt p)$. The terminal block can remain below
the full distance $\Theta(f)$, after adjusting constants and the minimum
size, because $\log f=O(f)$.

For a specified set of distinct bad resources, item 4 supplies disjoint
witnesses outside the large-cluster event. Choosing one witness and an
actual-fault subset for each resource gives disjoint specified sets of
actual locations. The same local-stochastic bound now applies to their
union, and the witness-counting sums factor as upper bounds even though
the noise is correlated. Thus the joint probability is bounded by
$q_{\rm inj}^{|J|}$. If the complement of the large-cluster event has
probability at least $1/2$, conditioning costs at most a factor two,
which is bounded by $2^{|J|}$ for nonempty $J$. Hence the conditional
resource family has $q_0\le2q_{\rm inj}=O(\sqrt p)$. A marginal
bound alone would not justify this conclusion.

\subsection{Distillation levels and total failure}
Writing $C_D=435$, the ideal recursion gives
\begin{equation}
 q_\ell\le C_D^{-1}(C_Dq_0)^{2^\ell}.
\end{equation}
Choose a fixed $p_{\rm th}$ small enough that $C_Dq_0<1$ uniformly
below threshold. Taking $2^\ell=\Theta(f)$ makes the resource error
exponentially small on the same scale as the protected-cluster error,
and $15^\ell=O(f^{\log_2 15})$. The number of locations and consumed
resources is at most $WT$ times a fixed polynomial in $f$ and
$15^\ell$. Its polynomial factors can be absorbed by increasing the
constant $c$ and minimum value $f_{\min}$ in Eq.~\eqref{eq:f}. A union
bound allocates at most $\varepsilon/2$ to large clusters and at most
$\varepsilon/2$ to final defective resources. This establishes the
claimed implementation-failure bound under the uniform hypotheses.

Code parameters give $K=\Theta(\sqrt f)$ and $n=\Theta(f^{3/2})$.
A coarse logical box of side $\Theta(f^{1/6})$ fits into one block with
fixed capacity reserved for live factory work. The number of blocks is
$O(W/K+B_f(A)+1)$, proving Eq.~\eqref{eq:space}. Depth-first evaluation
of a fixed-branching factory tree stores $O(\ell)$ candidates at once;
$\ell=O(\log f)$ fits within the reserved $\Theta(\sqrt f)$ capacity.
The geometrically increasing courier sizes have bounded total volume
relative to the largest one. Serial resource production then costs
$O(Kf15^\ell)$ per logical layer, proving Eq.~\eqref{eq:depth}.

For a bounded-aspect logical box, with $x=W/\sqrt f$,
$B_f(A)=O(x^{2/3}+1)$. Since $x^{2/3}\le x+1$, multiplication by
$f^{3/2}$ gives $O(Wf+f^{3/2})$. Connectedness alone is insufficient
for this bound: thin or fragmented regions can have many partially
occupied coarse boxes. The space and time conclusions use a serial
nearest-neighbor schedule with the specified hypotheses; they are not
an unconditional threshold proof for every code deformation sharing the
same static geometry.

\section{Numerical records and reproducibility}
\label{app:numerical}

The accompanying archive contains the revised \nolinkurl{manuscript.tex},
\nolinkurl{main_body.tex}, and \nolinkurl{appendices.tex}, plus a
\nolinkurl{research_records} directory copied from the supplied author
package. Paths in this appendix are relative to that directory.
The main numerical sources are as follows.
\begingroup\raggedright
\begin{itemize}
\item \nolinkurl{two_block_distance.csv} and
\nolinkurl{two_block_mc.csv}: initial interface distances, solver statuses,
and simulation counts. The implemented circuits and solver are in
\nolinkurl{two_block_measurement.py}.
\item \nolinkurl{numerical/scaling_campaign/mc_batches.csv}: unique
task/batch records, seeds, shot counts, any-observable failures, and
per-observable failures. The scheduler and profile are
\nolinkurl{scaling_campaign/campaign/run_mc.py} and
\nolinkurl{scaling_campaign/campaign/profiles.py}.
\item \nolinkurl{scaling_campaign/campaign/analyze.py} and
\nolinkurl{make_scaling_figure.py}: aggregation and suppression fits;
\nolinkurl{numerical/scaling_campaign/decoder_check.csv}: paired decoder
comparison; the same numerical directory contains code and fault-distance
results and the reported-bit witness checks.
\item \nolinkurl{numerical/continuous_transition/}: surface-code
transition campaign, circuit generator, counts, and posterior crossing
analysis. Headline runs use
\nolinkurl{final_transition_heavy.py}; posterior analysis is in
\nolinkurl{threshold_estimate/fit_threshold_estimate.py}.
\item \nolinkurl{growing_family/} and
\nolinkurl{routed_campaign_check.py}: routed-input campaign and its audit.
\item \nolinkurl{layer_code_explicit.py},
\nolinkurl{slab_gluing_certificate.py},
\nolinkurl{factory_map_audit.py}, and
\nolinkurl{boundary_composition_check.py}: finite construction, resource,
and boundary-map checks. \nolinkurl{schedule_verifier.py} checks a finite
coarse schedule; it does not prove a uniform microscopic detector catalog.
\end{itemize}
\endgroup

The revised numerical wording preserves archived observations and solver
statuses. No inference of an exact distance is made from a FEASIBLE
solver status alone. The independent revision checks reconstruct four
complete small layer-code cases, verify their ranks and commutation,
check 736 kernel-vector slab lifts at all positions, and check the
teleportation identities and scalar arithmetic. Their script and results
are supplied separately within the archive. These checks do not repeat
the full Monte Carlo campaigns or certify the architecture hypotheses.

\begin{table}[t]
\caption{Any-observable failures/shots for the spanning-grid suppression fits in Fig.~\ref{fig:standalone-scaling}(c). Counts aggregate unique archived batches. A dash denotes no data; zero-failure points are shown here but excluded by the archived fitting rule.}
\label{tab:fit-counts}
\centering\small
\setlength{\tabcolsep}{3pt}
\begin{tabular}{@{}lrrrrr@{}}
\toprule
Basis & $p$ ($10^{-3}$) & $s=2$ & $s=3$ & $s=4$ & $s=5$\\
\midrule
X & 0.5 & 52/68000 & 4/73200 & 0/25500 & 0/9900\\
X & 1 & 257/58000 & 29/40800 & 7/22200 & 0/9000\\
X & 2 & 49/2000 & 24/3000 & 25/6600 & 0/600\\
X & 3 & 135/2000 & 135/4800 & 45/2100 & 5/300\\
X & 4 & 289/2000 & 42/600 & -- & --\\
X & 6 & 802/2000 & 180/600 & 110/300 & --\\
\addlinespace
Z & 0.5 & 25/196000 & 1/72600 & 0/23400 & 0/9450\\
Z & 1 & 29/34000 & 18/69000 & 3/22800 & 1/9300\\
Z & 2 & 216/38000 & 31/14400 & 22/13500 & 1/1200\\
Z & 3 & 32/2000 & 55/4200 & 35/5100 & 0/150\\
Z & 4 & 90/2000 & 139/5400 & 65/2700 & --\\
Z & 6 & 322/2000 & 107/600 & 47/300 & --\\
\bottomrule
\end{tabular}
\end{table}

